\documentclass[11pt]{amsart}

\usepackage[T1]{fontenc}
\usepackage{lmodern}
\usepackage{amsmath,amssymb}
\usepackage{microtype}
\usepackage[
  colorlinks=true,
  linkcolor=blue,
  citecolor=blue,
  urlcolor=blue,
  pdftitle={DRAFT: Every K2,n is contractible},
  pdfauthor={Cavit Erginsoy}
]{hyperref}
\setlength{\emergencystretch}{1em}

\newtheorem{theorem}{Theorem}[section]
\newtheorem{lemma}[theorem]{Lemma}
\newtheorem{corollary}[theorem]{Corollary}

\title[DRAFT: Every $K_{2,n}$ is contractible]
{\textbf{DRAFT}\\[0.5ex]
Every $K_{2,n}$ is contractible}
\author{Cavit Erginsoy}
\date{Draft, 16 August 2026}

\subjclass[2020]{05C83, 05B35}
\keywords{graph schemes, rooted minors, contractible graphs, matroid union}

\begin{document}

\begin{abstract}
We prove that $K_{2,n}$ is contractible for every $n\geq1$: every graph
containing a $K_{2,n}$-scheme contains a rooted $K_{2,n}$ minor.  This
answers the $K_{2,4}$ part of Question~8.2 of K\"undgen, Pelsmajer and
Ramamurthi.  Since contractibility is inherited by subgraphs, every
bipartite graph with one part of order at most two is also contractible.
The proof encodes the two distinguished colour classes as connected
multigraphs on a common labelled edge set.  A component count and Edmonds'
matroid union theorem yield disjoint spanning trees.
\end{abstract}

\maketitle
\markboth{DRAFT: EVERY $K_{2,n}$ IS CONTRACTIBLE}
{CAVIT ERGINSOY: DRAFT}

\section{Introduction}

Let $G$ and $H$ be graphs with $V(H)\subseteq V(G)$.  An \emph{$H$-scheme}
in $G$ is a family
\[
                 (P_{uv}:uv\in E(H))
\]
of paths such that $P_{uv}$ is a $u$--$v$ path containing no other vertex
of $H$, and every non-empty collection of scheme paths with a common
vertex has a common end in $H$.  Their union is the \emph{underlying
graph} of the scheme.  A \emph{rooted $H$ minor} is a family
$(C_v:v\in V(H))$ of pairwise disjoint connected vertex sets such that
$C_u$ and $C_v$ are adjacent whenever $uv\in E(H)$, and $v\in C_v$ for
every $v\in V(H)$.  The graph $H$ is \emph{contractible} if every graph
containing an $H$-scheme contains a rooted $H$ minor.

K\"undgen, Pelsmajer and Ramamurthi introduced this terminology and
developed the basic theory of contractible graphs
\cite{KundgenPelsmajerRamamurthi2015}.  They proved that $K_{2,3}$ is
contractible \cite[Theorem~6.2]{KundgenPelsmajerRamamurthi2015}, and asked
whether $K_{2,4}$ or $K_{3,3}$ is contractible
\cite[Question~8.2]{KundgenPelsmajerRamamurthi2015}.  We answer the first
part and obtain the whole family at once.

\begin{theorem}\label{thm:main}
For every integer $n\geq1$, the graph $K_{2,n}$ is contractible.
\end{theorem}

The proof uses the coloured-scheme reduction of K\"undgen, Pelsmajer and
Ramamurthi and Edmonds' matroid union theorem.  The only additional
ingredient is the following elementary packing lemma.  Auxiliary
multigraphs may have parallel edges but no loops; all other graphs are
finite and simple.

\section{A two-projection packing lemma}

\begin{lemma}\label{lem:packing}
Let $m\geq1$, and let $M_1,M_2$ be connected multigraphs whose edge sets
are both labelled bijectively by the same set $E$.  For
$j\in\{1,2\}$, suppose that $M_j$ is the union of paths
$Q_1^j,\ldots,Q_m^j$ whose edge sets are pairwise disjoint and have union
$E$.  Trivial paths are allowed.  Suppose also that these paths all
contain a vertex $r_j$, and every other vertex of $M_j$ belongs to at
least two of them.  Then the graphic matroids of $M_1$ and $M_2$ have
disjoint bases.  Equivalently, there are disjoint subsets $T_1,T_2$ of
$E$ such that $T_j$ is a spanning tree of $M_j$.
\end{lemma}

\begin{proof}
For $X\subseteq E$, let $c_j(X)$ denote the number of components of the
spanning subgraph $(V(M_j),X)$, including isolated vertices, and put
$Y=E\setminus X$.  Count the incidences between these components and the
paths $Q_i^j$.  The component containing $r_j$ meets all $m$ paths, while
every other component meets at least two.  Hence
\[
        \sum_C |\{i:V(Q_i^j)\cap V(C)\ne\varnothing\}|
             \geq m+2(c_j(X)-1).                         \tag{2.1}
\]
Deleting the edges in $Y$ breaks $Q_i^j$ into at most
$|E(Q_i^j)\cap Y|+1$ pieces.  Distinct pieces may lie in the same
component of $(V(M_j),X)$, so the left side of (2.1) is at most
\[
             \sum_{i=1}^m (|E(Q_i^j)\cap Y|+1)=|Y|+m.
\]
It follows that
\[
                    c_j(X)\leq \frac{|Y|}{2}+1.           \tag{2.2}
\]

Let $\mathcal M_j$ be the graphic matroid of $M_j$, with rank function
$\rho_j$.  Since $M_j$ is connected,
\[
 \rho_j(E)=|V(M_j)|-1,
 \qquad
 \rho_j(X)=|V(M_j)|-c_j(X).
\]
Adding (2.2) for $j=1,2$ gives
\[
 |E\setminus X|+\rho_1(X)+\rho_2(X)
       \geq \rho_1(E)+\rho_2(E)                         \tag{2.3}
\]
for every $X\subseteq E$.  Edmonds' matroid union formula
\cite{Edmonds1970} says that
\[
 \rho_{\mathcal M_1\vee\mathcal M_2}(E)
   =\min_{X\subseteq E}
      \bigl(|E\setminus X|+\rho_1(X)+\rho_2(X)\bigr).
\]
By (2.3), this rank is at least $\rho_1(E)+\rho_2(E)$; taking $X=E$ in
the displayed minimum gives the reverse inequality.  Choose a maximum
independent set $I=I_1\cup I_2$ of the union, where $I_j$ is independent
in $\mathcal M_j$.  Its cardinality is $\rho_1(E)+\rho_2(E)$.  Equality in
\[
 |I_1\cup I_2|\leq |I_1|+|I_2|
                  \leq \rho_1(E)+\rho_2(E)
\]
shows that $I_1$ and $I_2$ are disjoint bases.  Taking $T_j=I_j$ proves
the lemma.
\end{proof}

\section{Proof of the theorem}

K\"undgen, Pelsmajer and Ramamurthi proved that an $H$-scheme can be
reduced, by root-preserving minor operations, to a \emph{coloured
$H$-scheme} \cite[Lemma~3.3]{KundgenPelsmajerRamamurthi2015}.  Such a
scheme has a proper colouring
\[
                         f:V(G)\longrightarrow V(H)
\]
which fixes $V(H)$, and $P_{uv}$ uses only the colours $u$ and $v$.
Moreover, every vertex outside $V(H)$ has degree at least four.  We shall
also use the following consequences recorded in
\cite[Remark~3.2(1), (2), and (7)]{KundgenPelsmajerRamamurthi2015}:
colours alternate on every scheme path; every edge lies on exactly one
scheme path; and, when $u$ has degree two in $H$, every vertex of colour
$u$ lies on both scheme paths incident with $u$.

\begin{proof}[Proof of Theorem~\ref{thm:main}]
Let the two vertices in the part of order two be $a,b$, and let the other
vertices be $t_1,\ldots,t_n$.  By the coloured-scheme reduction, it is
enough to find the required rooted minor in the underlying graph $G$ of a
coloured $K_{2,n}$-scheme.  Write
\[
 A=f^{-1}(a),\qquad B=f^{-1}(b),\qquad L_i=f^{-1}(t_i),
\]
and set
\[
                  E_i=L_i\setminus\{t_i\},
                  \qquad E=\bigcup_{i=1}^n E_i.          \tag{3.1}
\]

Every $x\in E_i$ lies internally on both $P_{at_i}$ and $P_{bt_i}$.
Thus $x$ has two distinct neighbours in $A$ on the first path and two
distinct neighbours in $B$ on the second.  Form an abstract
edge-labelled multigraph $M_a$ on vertex set $A$ by replacing each
$x\in E$ with an edge, labelled $x$, between its two neighbours in $A$.
Define $M_b$ on $B$ in the same way, using the same labelled edge set
$E$.

Suppressing the vertices of $E_i$ on $P_{at_i}$ gives a path $Q_i^a$ in
$M_a$, from $a$ to the $A$-neighbour of $t_i$.  If $P_{at_i}=at_i$, this
is the trivial path at $a$.  Alternation gives
\[
                         E(Q_i^a)=E_i.                    \tag{3.2}
\]
Thus the edge sets of $Q_1^a,\ldots,Q_n^a$ are pairwise disjoint and have
union $E$.  Every vertex of $A\setminus\{a\}$ lies on a scheme path
incident with $a$, so these paths also cover $V(M_a)$.  Their union is
therefore $M_a$, which is connected.  A vertex in $A\setminus\{a\}$
lying on $k$ of the paths has degree $2k$ in $G$: it is internal on each
path, and distinct scheme paths share no edge.  Since its degree is at
least four, $k\geq2$.  The same argument on the $b$ side gives paths
$Q_i^b$ with $E(Q_i^b)=E_i$, and every vertex of $B\setminus\{b\}$
belongs to at least two of them.

Lemma~\ref{lem:packing} now gives disjoint sets $T_a,T_b\subseteq E$ such
that $T_a$ is a spanning tree of $M_a$ and $T_b$ is a spanning tree of
$M_b$.  In $G$, put
\[
 \begin{aligned}
 C_a&=A\cup\{x\in E:x\in T_a\},\\
 C_b&=B\cup\{x\in E:x\in T_b\},\\
 C_i&=\{t_i\} \quad (1\leq i\leq n).
 \end{aligned}                                            \tag{3.3}
\]
These sets are pairwise disjoint.  An edge of $T_a$ labelled $x$, with
ends $u,v$, corresponds to the two-edge path with vertex sequence $u,x,v$
in $G$, so $G[C_a]$ is connected; similarly $G[C_b]$ is connected.  The
edge of $P_{at_i}$ incident with $t_i$ joins $C_i$ to $A\subseteq C_a$,
and the corresponding edge of $P_{bt_i}$ joins $C_i$ to $B\subseteq C_b$.
The sets in (3.3) therefore form a rooted $K_{2,n}$-minor model.

Finally, the coloured scheme lies in a rooted minor of the original
scheme graph.  Composing the two rooted minor models gives a rooted
$K_{2,n}$ minor in the original graph.
\end{proof}

Contractibility is inherited by subgraphs
\cite[Lemma~2.2]{KundgenPelsmajerRamamurthi2015}.  The theorem therefore
has the following immediate consequence.

\begin{corollary}
Every bipartite graph with a part of order at most two is contractible.
\end{corollary}

The argument is computation-free.  It does not address the $K_{3,3}$
part of Question~8.2 or the wider question of whether every bipartite
graph is contractible.

\section*{Statement on AI use and independence}

GPT-5.6 Sol, an OpenAI model, contributed to theorem discovery, proofs,
verification and drafting.  Cavit Erginsoy directed the work and made the
final decisions.  Erginsoy is not affiliated with OpenAI, and OpenAI
neither sponsored nor verified the work.

{\hbadness=2000
\bibliographystyle{amsplain}
\bibliography{references}
}

\end{document}
